% fixture: amsart/amsart-footnote
% tier: T3
% engine: document
% covers: amsart: its page and font sizes, \clubpenalty\@M, no \onecolumn \clearpage at the start, \thispagestyle{plain} as it loads, \@makefnmark (\raise.9ex at the \sf@size of the last math: amsart's \AtBeginDocument's \footnotesize, \LaTeXe's, a formula's, one in a footnote), \@makefntext and \footnoterule
% issue: #360
\documentclass{amsart}
\pagestyle{empty}
\begin{document}
Typesetting is the composition of text by means of arranging physical type or its digital equivalents. Stored letters and other symbols are retrieved and ordered according to the orthography of a language for visual display. The breaking of paragraphs into lines is a small part of this craft, but it decides how evenly the words are spread and how easily the eye moves from one line to the next.

Second paragraph with some more words to see indentation. Then a footnote\footnote{A note.} and more text.

Written in \LaTeXe\footnote{Its logo has a formula.} now.

After math $x+y$ the mark is at the script size\footnote{Of the normal size: $a+b$ here.} of the size of that math, and after the math in the note\footnote{Again.} at that of the footnote size.
\end{document}
